% Generated draft: requires mathematical and source review.
\documentclass[11pt]{article}
\usepackage[margin=25mm]{geometry}
\usepackage{fontspec}
\setmainfont{Latin Modern Roman}
\usepackage{amsmath,amssymb,mathtools,unicode-math}
\setmathfont{Latin Modern Math}
\usepackage{enumitem,needspace,xurl,hyperref}
\setlength{\parindent}{0pt}
\setlength{\parskip}{5pt}
\setlength{\emergencystretch}{3em}
\newcommand{\N}{\mathbb{N}}
\newcommand{\Z}{\mathbb{Z}}
\newcommand{\Q}{\mathbb{Q}}
\newcommand{\R}{\mathbb{R}}
\newcommand{\C}{\mathbb{C}}
\newcommand{\F}{\mathbb{F}}
\newcommand{\A}{\mathbb{A}}
\newcommand{\PP}{\mathbb{P}}
\newcommand{\E}{\mathbb{E}}
\newcommand{\eps}{\varepsilon}
\newcommand{\dd}{\,\mathrm{d}}
\begin{document}
\section*{UnsolvedMath ID: 30004186}
\textbf{Source problem code:} catalogue entry
\subsection{Definitions and mathematical statement}
A traveling wave has the form \(u(x,t)=U_*(z)\), where \(z=x-ct\) and \(c\) is its propagation speed. The supplied profile formula is
\[
U_*(z)=\frac{3z^2-\pi^2}{18},
\]
with periodic extension understood. To specify that extension, one must prescribe a fundamental interval \(I\) of length \(L>0\), use the displayed formula on \(I\), and impose \(U_*(z+L)=U_*(z)\). The supplied record does not explicitly specify \(I\), \(L\), or \(c\); thus this description does not select a unique periodic profile.

\textbf{Source question.} For the reduced Ostrovsky equation, determine whether the indicated peaked periodic traveling waves are nonlinearly stable or nonlinearly unstable.

\textbf{Unresolved qualifications.} The record does not give the equation's normalization, the spatial domain, the solution concept, the admissible perturbations, the function space or topology, or the time interval over which stability is required. It also does not specify whether stability is measured relative to a fixed wave or modulo spatial translation. These choices determine the hypotheses and quantifiers of a nonlinear stability statement, so a unique precise stability assertion cannot be extracted from the supplied record alone.
\subsection{Short English statement}
Determine the nonlinear stability or instability of the periodically extended peaked traveling wave with profile \(U_*(z)=(3z^2-\pi^2)/18\) for the reduced Ostrovsky equation. The periodic extension, equation normalization, and meaning of nonlinear stability require clarification before the question has a fully specified mathematical formulation.
\subsection{Sources}
\begin{itemize}
\item Supplied local UnsolvedMath registry record, permanent ID \(30004186\), source problem code \texttt{catalogue entry}, titled \emph{Stability of Peaked Reduced Ostrovsky Waves}, in the Oberwolfach Reports Open Problems collection. The proposed year is recorded as 2019; no individual proposer is supplied.
\item Supplied report citation: \emph{Mathematical Theory of Water Waves (2020)}, Oberwolfach Reports, pp.~1919--1979, DOI \texttt{10.4171/owr/2019/32}. \url{https://doi.org/10.4171/owr/2019/32}. The background describes these reports as documenting mathematical workshops and their open-problem discussions.
\item Supplied literature reference: Geyer--Pelinovsky, linear instability. \url{https://doi.org/10.1137/18M117978X}.
\item Supplied literature reference: Geyer--Pelinovsky, spectral instability. \url{https://doi.org/10.1090/proc/14937}.
\item Supplied literature reference: original OWR report. \url{https://doi.org/10.4171/owr/2019/32}.
\item Literature assessment checked at \texttt{2026-08-22}: the supplied curation labels the problem partially solved and reports that linear and spectral instability, together with uniqueness, are proved for the peaked periodic wave, while nonlinear instability and the full nonlinear stability question remain open. This is supplied dated curation derived from the OWR collection, not an independently verified assessment; its time-sensitive status and cited claims require independent verification.
\end{itemize}
\end{document}
